% !TeX program = lualatex
\documentclass[12pt,aspectratio=169]{beamer}
\usepackage[utf8]{inputenc}
\usepackage[T1]{fontenc}
\usepackage{amsmath,amsfonts,amssymb}
\usepackage{graphicx}
\usepackage{tikz}
\usepackage{xcolor}
\usepackage{hyperref}
\usepackage{anyfontsize}
\usepackage{lmodern}

% ====== MAKE ALL FONTS SMALLER (GLOBAL SCALING) ======
\renewcommand{\tiny}{\fontsize{5}{6}\selectfont}
\renewcommand{\scriptsize}{\fontsize{6}{7}\selectfont}
\renewcommand{\footnotesize}{\fontsize{7}{8}\selectfont}
\renewcommand{\small}{\fontsize{8}{9}\selectfont}
\renewcommand{\normalsize}{\fontsize{9}{10}\selectfont}
\renewcommand{\large}{\fontsize{10}{11}\selectfont}
\renewcommand{\Large}{\fontsize{11}{13}\selectfont}
\renewcommand{\LARGE}{\fontsize{13}{15}\selectfont}
\renewcommand{\huge}{\fontsize{15}{17}\selectfont}
\renewcommand{\Huge}{\fontsize{18}{21}\selectfont}

% ====== COLOR THEME ======
\definecolor{redcorrect}{RGB}{200,30,30}
\definecolor{bluecorrect}{RGB}{30,80,200}
\definecolor{greencheck}{RGB}{0,150,50}
\definecolor{redx}{RGB}{200,0,0}
\definecolor{lightgray}{RGB}{245,245,245}
\definecolor{darkgray}{RGB}{50,50,50}
\definecolor{softyellow}{RGB}{255,248,200}
\definecolor{infoblue}{RGB}{230,240,252}

\setbeamercolor{title}{fg=white,bg=redcorrect}
\setbeamercolor{frametitle}{fg=white,bg=redcorrect}
\setbeamercolor{block title}{fg=white,bg=bluecorrect}
\setbeamercolor{block body}{fg=darkgray,bg=lightgray}
\setbeamercolor{normal text}{fg=darkgray}
\usecolortheme{default}

% ====== CUSTOM COMMANDS ======
\newcommand{\correct}[1]{\textcolor{greencheck}{\textbf{\checkmark}} \textcolor{greencheck}{#1}}
\newcommand{\wrong}[1]{\textcolor{redx}{\textbf{\times}} \textcolor{redx}{#1}}
\newcommand{\highlightred}[1]{\textcolor{redcorrect}{\textbf{#1}}}
\newcommand{\highlightblue}[1]{\textcolor{bluecorrect}{\textbf{#1}}}
\newcommand{\tipbox}[1]{\fcolorbox{bluecorrect}{softyellow}{\parbox{0.88\textwidth}{\centering #1}}}
\newcommand{\wrongbox}[1]{\fcolorbox{redx}{red!10}{\parbox{0.88\textwidth}{\wrong{#1}}}}

% ====== NO ZOOM — JUST SMALLER FONTS EVERYWHERE ======
% Each slide uses \small or \footnotesize to shrink content

\begin{document}
	
	% ================================================================
	% SLIDE 1: TITLE
	% ================================================================
	\begin{frame}
		\centering
		\vspace{0.5cm}
		{\fontsize{18}{22}\selectfont \textbf{Money Laundering and Financial Crime}}\\[0.2cm]
		{\fontsize{14}{17}\selectfont Test Review \& Answer Explanations}\\[0.4cm]
		{\fontsize{11}{14}\selectfont Compliance Training Department}\\[0.1cm]
		{\fontsize{9}{11}\selectfont \today}
		\vspace{0.3cm}
		\rule{4cm}{1pt}
	\end{frame}
	
	% ================================================================
	% SLIDE 2: WHY SO HARD?
	% ================================================================
	\begin{frame}
		\frametitle{\fontsize{13}{16}\selectfont Why Is It So Hard to Fight Financial Crime Globally?}
		
		\small
		\begin{block}{\fontsize{11}{13}\selectfont The Problem: Different Countries, Different Rules}
			\footnotesize
			\begin{itemize}
				\item Imagine playing a sport where \textbf{every country has different rules}
				\item One country says an action is \textbf{legal}, another says it is \textbf{criminal}
				\item Criminals \textbf{exploit these gaps} — they shift activities to the weakest link
			\end{itemize}
		\end{block}
		
		\vspace{0.05cm}
		
		\begin{columns}[T]
			\column{0.48\textwidth}
			\footnotesize
			\highlightred{\textbf{What I Got Wrong:}}
			\begin{itemize}
				\item I thought \textit{digital networks} were the problem
				\item But digital networks are \textbf{not} obsolete — they are \textbf{getting better}
				\item AI, machine learning, and blockchain analytics \textbf{help} detect crime
			\end{itemize}
			
			\column{0.48\textwidth}
			\footnotesize
			\highlightblue{\textbf{What I Should Have Focused On:}}
			\begin{itemize}
				\item The real challenge is \textbf{legal divergence}
				\item Countries \textbf{criminalize different behaviors}
				\item Example: "Bribery" in one country may be "normal hospitality" in another
				\item This creates \textbf{gaps} that are hard to close
			\end{itemize}
		\end{columns}
		
		\vspace{0.15cm}
		\tipbox{\footnotesize \textbf{Core Lesson:} The problem is not technology — it is \textit{human} systems. Legal fragmentation is the real enemy.}
	\end{frame}
	
	% ================================================================
	% SLIDE 3: CULTURAL HOSPITALITY
	% ================================================================
	\begin{frame}
		\frametitle{\fontsize{13}{16}\selectfont Cultural Hospitality — When "Being Nice" Creates Risk}
		
		\small
		\begin{block}{\fontsize{11}{13}\selectfont The Scenario}
			\footnotesize
			A consultant reviews a company's anti-bribery policy and notices something missing: \\
			\medskip
			\centering
			\textit{"Lavish hospitality extended to visiting foreign officials is not clearly addressed."}
		\end{block}
		
		\vspace{0.05cm}
		
		\footnotesize
		\highlightred{\textbf{Why This Is a Problem:}}
		\begin{itemize}
			\item In some cultures, \textbf{expensive gifts} and \textbf{lavish entertainment} are expected
			\item But \textbf{where is the line?} When does hospitality become bribery?
			\item Without clear guidance, employees \textbf{may cross that line without realizing it}
		\end{itemize}
		
		\vspace{0.03cm}
		
		\highlightblue{\textbf{The Risk This Creates:}}
		\begin{itemize}
			\item The organization could \textbf{inadvertently facilitate bribery}
			\item It happens \textit{under the guise} of "cultural expectations"
			\item The company did not \textit{intend} to bribe — but the action still constitutes bribery
		\end{itemize}
		
		\vspace{0.1cm}
		\tipbox{\footnotesize \textbf{Best Practice:} Set clear monetary limits, require pre-approval, and train staff to distinguish \textit{culture} from \textit{corruption}.}
	\end{frame}
	
	% ================================================================
	% SLIDE 4: THREE STAGES
	% ================================================================
	\begin{frame}
		\frametitle{\fontsize{13}{16}\selectfont The Three Stages — Where Do Shell Companies Fit?}
		
		\small
		\begin{block}{\fontsize{11}{13}\selectfont A Quick Refresher: The Money Laundering Stages}
			\footnotesize
			\begin{enumerate}
				\item \textbf{Placement:} Get dirty cash into the system (deposits, currency exchange)
				\item \textbf{Layering:} Move it around to hide where it came from (multiple transactions, shell companies)
				\item \textbf{Integration:} Use it as "clean" money (investments, purchases)
			\end{enumerate}
		\end{block}
		
		\vspace{0.05cm}
		
		\begin{columns}[T]
			\column{0.48\textwidth}
			\footnotesize
			\highlightred{\textbf{What I Got Wrong:}}
			\begin{itemize}
				\item I thought shell companies are used in \textbf{placement}
				\item But placement is about \textbf{getting cash in}
				\item Shell companies are not needed for that — cash is too visible
			\end{itemize}
			
			\column{0.48\textwidth}
			\footnotesize
			\highlightblue{\textbf{The Correct Answer:}}
			\begin{itemize}
				\item Shell companies are used in \textbf{layering}
				\item They create \textbf{complex chains} of transactions
				\item They \textbf{obscure beneficial ownership}
				\item They make money look like it is from \textit{legitimate business}
			\end{itemize}
		\end{columns}
		
		\vspace{0.15cm}
		\tipbox{\footnotesize \textbf{Memory Trick:} \textit{"Placement = Put it in; Layering = Hide it; Integration = Use it."} Shell companies = \textbf{HIDING} = Layering.}
	\end{frame}
	
	% ================================================================
	% SLIDE 5: MOBILE ONBOARDING
	% ================================================================
	\begin{frame}
		\frametitle{\fontsize{13}{16}\selectfont Mobile Onboarding — Balancing Compliance and Customer Access}
		
		\small
		\begin{block}{\fontsize{11}{13}\selectfont The Challenge}
			\footnotesize
			A building society wants to offer services through its mobile app. \\
			But they are worried about \textbf{identity fraud}. \\
			\medskip
			\centering
			\textit{"How do we stay compliant without ruining the customer experience?"}
		\end{block}
		
		\vspace{0.05cm}
		
		\footnotesize
		\highlightred{\textbf{Why This Is Tricky:}}
		\begin{itemize}
			\item \textbf{Compliance} wants \textit{strong} identity verification
			\item \textbf{Customers} want \textit{fast} and \textit{easy} access
			\item These two goals often \textbf{pull in opposite directions}
		\end{itemize}
		
		\vspace{0.03cm}
		
		\highlightblue{\textbf{The Best Solution:}}
		\begin{itemize}
			\item \textbf{Biometric verification} (fingerprint, facial recognition) — strong and easy
			\item \textbf{Digital onboarding review} — ensure the \textit{process itself} is secure
			\item This gives you \textbf{both}: strong identity \textbf{and} a smooth experience
		\end{itemize}
		
		\vspace{0.1cm}
		\tipbox{\footnotesize \textbf{Key Takeaway:} Don't choose between compliance and access — use \textit{technology} to achieve both.}
	\end{frame}
	
	% ================================================================
	% SLIDE 6: DEFENSIVE SARS
	% ================================================================
	\begin{frame}
		\frametitle{\fontsize{13}{16}\selectfont Defensive SARs — Why "Just in Case" Is a Problem}
		
		\small
		\begin{block}{\fontsize{11}{13}\selectfont What Is a Defensive SAR?}
			\footnotesize
			It is a Suspicious Activity Report filed \textbf{not because there is genuine suspicion}, but because the filer wants to \textbf{avoid blame} later.
		\end{block}
		
		\vspace{0.05cm}
		
		\footnotesize
		\highlightred{\textbf{What I Got Wrong:}}
		\begin{itemize}
			\item I thought defensive SARs \textbf{simplify} compliance
			\item The truth is \textbf{they do the opposite}
			\item They create \textbf{more work}, not less
		\end{itemize}
		
		\vspace{0.03cm}
		
		\highlightblue{\textbf{The Real Problem:}}
		\begin{itemize}
			\item Regulators receive \textbf{thousands} of SARs every day
			\item Defensive filings \textbf{flood} the system with \textit{irrelevant} reports
			\item This creates \textbf{noise} — making it harder to find \textbf{real} suspicious activity
			\item It \textbf{burdens regulators} with excessive and potentially irrelevant reports
		\end{itemize}
		
		\vspace{0.1cm}
		\tipbox{\footnotesize \textbf{Rule to Remember:} File SARs only when you have \textit{reasonable suspicion} — not to "cover your back."}
	\end{frame}
	
	% ================================================================
	% SLIDE 7: LAW ENFORCEMENT CALL
	% ================================================================
	\begin{frame}
		\frametitle{\fontsize{13}{16}\selectfont Responding to Law Enforcement — The Safe Approach}
		
		\small
		\begin{block}{\fontsize{11}{13}\selectfont The Scenario}
			\footnotesize
			A law enforcement officer calls a financial crime investigator. \\
			The officer wants details about a suspicious customer transaction. \\
			\medskip
			\centering
			\textit{The call is informal, with no written documentation.}
		\end{block}
		
		\vspace{0.05cm}
		
		\footnotesize
		\highlightred{\textbf{Why This Situation Is Risky:}}
		\begin{itemize}
			\item Without written documentation, there is \textbf{no official record}
			\item The investigator could \textbf{misremember} details
			\item The call could be \textbf{impersonation} — a fraudster pretending to be law enforcement
		\end{itemize}
		
		\vspace{0.03cm}
		
		\highlightblue{\textbf{The Best Next Step:}}
		\begin{itemize}
			\item \textbf{Document the key points} of the conversation immediately
			\item \textbf{Request formal written confirmation} from the officer
			\item This creates an \textbf{audit trail}
			\item It protects both the investigator \textbf{and} the institution
		\end{itemize}
		
		\vspace{0.1cm}
		\tipbox{\footnotesize \textbf{Always Remember:} If it is not documented, it did not happen. \textit{Always} get it in writing.}
	\end{frame}
	
	% ================================================================
	% SLIDE 8: CONCENTRATION ACCOUNTS
	% ================================================================
	\begin{frame}
		\frametitle{\fontsize{13}{16}\selectfont Concentration Accounts — When Oversight Is Missing}
		
		\small
		\begin{block}{\fontsize{11}{13}\selectfont What Is a Concentration Account?}
			\footnotesize
			An internal account used to \textbf{aggregate} multiple transactions before they are posted to individual customer accounts.
		\end{block}
		
		\vspace{0.05cm}
		
		\footnotesize
		\highlightred{\textbf{The Regulatory Concern:}}
		\begin{itemize}
			\item If a concentration account lacks \textbf{sufficient oversight}, it becomes a \textbf{black box}
			\item Funds go \textit{in} and come \textit{out} — but it is \textbf{hard to trace} where they went
			\item Criminals can use this \textbf{poor traceability} to move money undetected
		\end{itemize}
		
		\vspace{0.03cm}
		
		\highlightblue{\textbf{Why This Matters:}}
		\begin{itemize}
			\item The institution might \textbf{inadvertently facilitate money laundering}
			\item Not because they \textit{want} to, but because the account \textbf{lacks proper controls}
			\item This is why regulators \textbf{scrutinize} concentration accounts closely
		\end{itemize}
		
		\vspace{0.1cm}
		\tipbox{\footnotesize \textbf{Core Principle:} If you cannot \textit{trace} it, you cannot \textit{control} it. Oversight is everything.}
	\end{frame}
	
	% ================================================================
	% SLIDE 9: PRIORITIZING REVIEWS
	% ================================================================
	\begin{frame}
		\frametitle{\fontsize{13}{16}\selectfont Prioritizing Reviews — Not All Trades Are Equal}
		
		\small
		\begin{block}{\fontsize{11}{13}\selectfont The Scenario}
			\footnotesize
			A compliance officer evaluates two trades: \\
			\medskip
			\begin{itemize}
				\item \textbf{Trade A:} High-volume equities with normal volatility
				\item \textbf{Trade B:} Low-volume penny stocks with rapid pricing shifts
			\end{itemize}
			Both passed basic monitoring. Which should be prioritized?
		\end{block}
		
		\vspace{0.05cm}
		
		\footnotesize
		\highlightred{\textbf{What I Got Wrong:}}
		\begin{itemize}
			\item I thought both should be treated \textbf{equally}
			\item I assumed volume and volatility are \textbf{neutral indicators}
			\item This is a \textbf{dangerous assumption}
		\end{itemize}
		
		\vspace{0.03cm}
		
		\highlightblue{\textbf{The Correct Priority:}}
		\begin{itemize}
			\item \textbf{The penny stock trade} should be prioritized
			\item Why? Because penny stocks have \textbf{higher manipulation risk}
			\item They are easier to \textbf{manipulate} — small trades can move the price
			\item Criminals use them to \textbf{create artificial activity}
		\end{itemize}
		
		\vspace{0.1cm}
		\tipbox{\footnotesize \textbf{Key Insight:} Do not assume all activity is equal. \textit{Lower volume} + \textit{higher volatility} = \textbf{higher risk}.}
	\end{frame}
	
	% ================================================================
	% SLIDE 10: FACIAL RECOGNITION
	% ================================================================
	\begin{frame}
		\frametitle{\fontsize{13}{16}\selectfont Facial Recognition — Why Biometrics Come With a Catch}
		
		\small
		\begin{block}{\fontsize{11}{13}\selectfont The Comparison}
			\footnotesize
			\begin{itemize}
				\item \textbf{Passwords:} Can be \textit{changed} if compromised
				\item \textbf{Biometrics (facial recognition):} \textbf{Cannot} be changed easily
			\end{itemize}
		\end{block}
		
		\vspace{0.05cm}
		
		\footnotesize
		\highlightred{\textbf{What I Got Wrong:}}
		\begin{itemize}
			\item I thought the problem was \textbf{storage capacity}
			\item Storage is \textbf{not} the main issue with biometrics
		\end{itemize}
		
		\vspace{0.03cm}
		
		\highlightblue{\textbf{The Real Risk:}}
		\begin{itemize}
			\item If a password is stolen, you \textbf{reset it} — simple
			\item If your \textbf{biometric data} (face, fingerprint) is compromised, you \textbf{cannot reset it}
			\item You have \textit{one} face, \textit{one} set of fingerprints
			\item Once stolen, that data is \textbf{permanently} exposed
			\item This is the \textbf{fundamental vulnerability} of biometric authentication
		\end{itemize}
		
		\vspace{0.1cm}
		\tipbox{\footnotesize \textbf{Critical Lesson:} Passwords are replaceable. Biometrics are \textit{not}. That is why they must be \textbf{protected even more carefully}.}
	\end{frame}
	
	% ================================================================
	% SLIDE 11: AML CULTURE
	% ================================================================
	\begin{frame}
		\frametitle{\fontsize{13}{16}\selectfont AML Culture — When Policies Are Not Enough}
		
		\small
		\begin{block}{\fontsize{11}{13}\selectfont The Finding}
			\footnotesize
			An internal audit finds something concerning: \\
			\medskip
			\centering
			\textit{"Despite robust policies, employee engagement on compliance is low. Staff view compliance as a barrier to performance."}
		\end{block}
		
		\vspace{0.05cm}
		
		\footnotesize
		\highlightred{\textbf{What I Got Wrong:}}
		\begin{itemize}
			\item I assumed the policies were \textbf{too complex}
			\item Complexity is \textbf{not} the root issue here
		\end{itemize}
		
		\vspace{0.03cm}
		
		\highlightblue{\textbf{The Real Problem:}}
		\begin{itemize}
			\item The organization has \textbf{failed to embed compliance values} in its operational culture
			\item Having \textbf{policies on paper} is not enough
			\item Employees must \textbf{believe} in compliance — not just follow rules
			\item When compliance is seen as a \textbf{barrier}, people will \textbf{find ways around it}
		\end{itemize}
		
		\vspace{0.1cm}
		\tipbox{\footnotesize \textbf{Core Lesson:} Compliance is not a set of rules — it is a \textit{culture}. If you do not build the culture, the rules will not work.}
	\end{frame}
	
	% ================================================================
	% SLIDE 12: BOARD OVERSIGHT
	% ================================================================
	\begin{frame}
		\frametitle{\fontsize{13}{16}\selectfont Board Oversight — The Director's Role in AML}
		
		\small
		\begin{block}{\fontsize{11}{13}\selectfont Effective Board Oversight Includes:}
			\footnotesize
			\begin{enumerate}
				\item \textbf{Receiving regular reports} on financial crime risks and control effectiveness
				\item \textbf{Reviewing whether senior management} is addressing audit findings
				\item \textbf{Challenging management} when AML performance metrics show signs of weakness
			\end{enumerate}
		\end{block}
		
		\vspace{0.05cm}
		
		\footnotesize
		\highlightred{\textbf{What I Got Wrong:}}
		\begin{itemize}
			\item I thought board members should \textbf{participate in SAR investigations}
			\item This is \textbf{operational work} — boards should not do this
			\item Board members are \textbf{oversight}, not \textit{operators}
		\end{itemize}
		
		\vspace{0.03cm}
		
		\highlightblue{\textbf{The Correct Role of the Board:}}
		\begin{itemize}
			\item \textbf{Oversee} — do not \textit{do}
			\item Ask tough questions — \textbf{challenge} management
			\item Ensure \textbf{accountability} — is management acting on audit findings?
			\item Monitor \textbf{trends} — are risks increasing or decreasing?
		\end{itemize}
		
		\vspace{0.1cm}
		\tipbox{\footnotesize \textbf{Board Principle:} Directors oversee, they do not operate. Their job is to \textit{challenge} and \textit{question}.}
	\end{frame}
	
	% ================================================================
	% SLIDE 13: TERRORIST FINANCING
	% ================================================================
	\begin{frame}
		\frametitle{\fontsize{13}{16}\selectfont Terrorist Financing — How Money Moves Across Borders}
		
		\small
		\begin{block}{\fontsize{11}{13}\selectfont Common Techniques Used by Terrorists:}
			\footnotesize
			\begin{enumerate}
				\item \textbf{Informal value transfer systems} (Hawala, Hundi)
				\item \textbf{Physical transportation of cash} across porous borders
				\item \textbf{Fake charities} operating in high-risk regions
			\end{enumerate}
		\end{block}
		
		\vspace{0.05cm}
		
		\footnotesize
		\highlightred{\textbf{What I Missed:}}
		\begin{itemize}
			\item I selected only \textit{two} when the question asked for \textit{three}
			\item I missed \textbf{physical cash transportation}
			\item Terrorists often simply \textbf{carry cash} across borders
			\item This is old-school but \textbf{still effective}
		\end{itemize}
		
		\vspace{0.03cm}
		
		\highlightblue{\textbf{Why These Three Are Correct:}}
		\begin{itemize}
			\item \textbf{Hawala:} Trust-based, no paper trail, hard to detect
			\item \textbf{Cash:} Simple, no system to flag, works in vulnerable regions
			\item \textbf{Charities:} Appear legitimate, move money under humanitarian cover
		\end{itemize}
		
		\vspace{0.1cm}
		\tipbox{\footnotesize \textbf{Remember:} Terrorists use both \textit{high-tech} and \textit{no-tech} methods. Do not overlook the simple ones.}
	\end{frame}
	
	% ================================================================
	% SLIDE 14: HUMAN TRAFFICKING
	% ================================================================
	\begin{frame}
		\frametitle{\fontsize{13}{16}\selectfont Human Trafficking — Laundering the Proceeds}
		
		\small
		\begin{block}{\fontsize{11}{13}\selectfont Characteristics That Enable Laundering:}
			\footnotesize
			\begin{enumerate}
				\item \textbf{Use of mobile payment and remittance systems}
				\item \textbf{Weak cross-border enforcement cooperation}
			\end{enumerate}
		\end{block}
		
		\vspace{0.05cm}
		
		\footnotesize
		\highlightred{\textbf{What I Got Wrong:}}
		\begin{itemize}
			\item I thought counterfeit documents were a key factor
			\item While documents \textit{are} used in trafficking, they are not the \textbf{primary laundering enabler}
			\item I also missed \textbf{mobile payments} as a key technique
		\end{itemize}
		
		\vspace{0.03cm}
		
		\highlightblue{\textbf{Why These Two Are Correct:}}
		\begin{itemize}
			\item \textbf{Mobile payments} are hard to trace, fast, and cross-border
			\item Traffickers use them to \textbf{move small amounts} frequently
			\item \textbf{Weak enforcement} means criminals can operate across borders without fear
		\end{itemize}
		
		\vspace{0.1cm}
		\tipbox{\footnotesize \textbf{Key Insight:} Traffickers use \textit{accessible} systems (mobile money) and exploit \textit{jurisdictional gaps} (weak cooperation).}
	\end{frame}
	
	% ================================================================
	% SLIDE 15: FINAL SUMMARY
	% ================================================================
	\begin{frame}
		\frametitle{\fontsize{13}{16}\selectfont Final Summary — What I Learned From This Review}
		
		\small
		\begin{block}{\fontsize{11}{13}\selectfont The Most Important Takeaways:}
			\vspace{0.03cm}
			\footnotesize
			\begin{enumerate}
				\item \textbf{Legal divergence} is the real challenge in global AML — not technology
				\item \textbf{Cultural practices} can blur the line between hospitality and bribery
				\item Shell companies belong in \textbf{layering}, not placement
				\item \textbf{Biometrics} are powerful but \textit{cannot be reset} — protect them
				\item \textbf{Defensive SARs} flood the system and make detection harder
				\item \textbf{Document everything} — especially law enforcement requests
				\item \textbf{Concentration accounts} need strong oversight to prevent misuse
				\item Not all trades are equal — \textbf{prioritize} based on risk, not volume
				\item \textbf{Compliance culture} matters more than compliance policies
				\item Boards \textbf{oversee} — they do not \textit{operate}
			\end{enumerate}
		\end{block}
		
		\vspace{0.08cm}
		\centering
		\fcolorbox{redcorrect}{softyellow}{\parbox{0.85\textwidth}{\centering
				\footnotesize \textbf{Remember:} Every wrong answer is a learning opportunity. \\
				\textit{Understand why — and you will not make the same mistake again.}}}
	\end{frame}
	
	% ================================================================
	% SLIDE 16: THANK YOU
	% ================================================================
	\begin{frame}
		\centering
		\vspace{2.5cm}
		{\fontsize{22}{26}\selectfont \textbf{Thank You}} \\[0.3cm]
		{\fontsize{14}{17}\selectfont Keep learning. Keep growing.} \\[0.2cm]
		\rule{3cm}{1pt} \\[0.2cm]
		{\fontsize{10}{12}\selectfont \textit{Compliance is not about perfection — it is about continuous improvement.}}
	\end{frame}
	
\end{document}